% Propietario conservado: /Users/ruben/Documents/New project/output/EXTREMA_MEDIA_RAZON_RECIPROCIDAD_ES_EN_20260918/ARTICULO/ES/source/libro/reciprocidad_radial.tex
\chapter{Géométrie hyperbolique de la paire angulaire et récupération de l'action}
\label{x_reciprocidad:rec:radial}

\section{Radion, ellipse d’action et cartes marquées}

La section orientée de phase, d’action et de mémoire dénommée \emph{radion} conserve conjointement l’incrément angulaire \(\theta=1\), l’action signée \(\sigma\hbar_{\rm ret}\), l’appartenance au tour \(2\pi_{\HMT}\), le feuillet, le retour nonadique et le cylindre de raffinement. Son raccordement d’unité s’exprime dans la carte du continu par \(1=S_E-\Phi_T+\varepsilon_1\). La projection qui retient seulement \(\theta=1\) produit le radian ordinaire ; la section enrichie conserve en outre le registre nécessaire pour transporter et inverser ses autres lectures. Dans cette section, \emph{rayon} désignera une coordonnée positive de forme, distincte de cet objet enrichi.


Pour expliciter les termes du raccord, la roue dodécaphasique a pour centres \(30m-10\) dans sa carte en degrés. La seconde phase apporte \(50\), et la publication \(A=1000\alpha\) détermine
\[
 S_E=\frac{2\pi}{360}(50+1000\alpha).
\]
Le secteur actif du sélecteur tritique réunit \(H_+=\{1,4,7\}\) et \(H_-=\{2,5,8\}\). Les quinze paires non ordonnées, lues dans les six positions actives, produisent \(N_6=6\binom62=90\) ; leur extension à huit positions produit \(N_8=8\binom62=120\). Les deux feuillets du secteur actif, normalisés par la fibre ambiante de \(729\) éléments, fixent
\[
 \rho_{\rm tw}=\frac{2N_6}{729}=\frac{20}{81},\qquad
 \Phi_T=\frac{20}{81}\Delta_4,\qquad
 S_T=1+\Phi_T,\qquad \varepsilon_1=S_T-S_E,
\]
où \(\Delta_4\) est la coordonnée torsionnelle de \eqref{x_reciprocidad:rec:eq:electron-basal}. La soustraction conserve les sections qui la précèdent et prouve directement
\(S_E-\Phi_T+\varepsilon_1=S_E-\Phi_T+1+\Phi_T-S_E=1\).
Le coefficient provient du recensement et des deux feuillets, et non de la valeur de la différence finale.

Les publications angulaires des sections précédentes satisfont, dans la carte considérée,
\[
 A>0,\qquad |C^*|<A,\qquad
 \xi=\frac{C^*}{A},\qquad \eta_{\rm el}=\operatorname{artanh}\xi.
\]
La notation \(\eta_{\rm el}\) distingue l'anisotropie elliptique du coefficient de retour \(\eta_{\rm ret}\). Pour une section positive d'action \(\hbar\), les demi-axes sont
\begin{equation}
 a_S=\sqrt{2\hbar}\,e^{\eta_{\rm el}/2},\qquad
 b_S=\sqrt{2\hbar}\,e^{-\eta_{\rm el}/2},\qquad a_Sb_S=2\hbar.
 \label{x_reciprocidad:rec:eq:elipse}
\end{equation}
Le produit fixe l’action par tour, tandis que le quotient fixe la forme. Nous définissons
\[
 R_\star=\sqrt{b_S/a_S}=e^{-\eta_{\rm el}/2},\qquad
 \tau_{\rm el}=iR_\star^2.
\]

\begin{theorem}[Compatibilité des cartes radiale, de Poincaré et de Klein]
\label{x_reciprocidad:rec:thm:klein}
Soient
\[
 w=\frac{\tau_{\rm el}-i}{\tau_{\rm el}+i},\qquad
 x_K=\frac{2w}{1+w^2}.
\]
Dans le diamètre réel du disque de Poincaré et son image dans celui de Klein,
\begin{align}
 w&=\frac{R_\star^2-1}{R_\star^2+1}
      =-\tanh\frac{\eta_{\rm el}}2,\label{x_reciprocidad:rec:eq:poincare}\\
 x_K&=\frac{R_\star^4-1}{R_\star^4+1}
      =-\tanh\eta_{\rm el}=-\frac{C^*}{A}.\label{x_reciprocidad:rec:eq:klein}
\end{align}
Toutes ces transformations sont inversibles dans leurs domaines ; en particulier,
\(R_\star^2=(1+w)/(1-w)\).
\end{theorem}
\begin{proof}
On substitue \(\tau=iR_\star^2\) dans la transformation fractionnaire. Le quotient simplifié a pour dénominateur \(R_\star^2+1>0\) et \(-1<w<1\). En substituant ce quotient dans \(2w/(1+w^2)\) et en multipliant le numérateur et le dénominateur par \((R_\star^2+1)^2\), on obtient \((R_\star^4-1)/(R_\star^4+1)\). L’égalité \(R_\star^4=e^{-2\eta_{\rm el}}\) donne la tangente hyperbolique. Enfin, \(\tanh\eta_{\rm el}=\xi\). Les récupérateurs s’obtiennent en isolant les quotients.
\end{proof}

Le signe de \eqref{x_reciprocidad:rec:eq:klein} correspond à l'orientation choisie pour le diamètre. L'orientation opposée change les deux signes. C'est une coordonnée hyperbolique de la carte rectangulaire marquée : elle conserve la marque et ne se confond ni avec une forme modulaire ni avec une coordonnée du quotient modulaire dépourvu d'orientation.

L’échange des demi-axes agit par
\begin{equation}
 (R_\star,\eta_{\rm el},\tau_{\rm el},w,x_K)
 \longmapsto(R_\star^{-1},-\eta_{\rm el},-1/\tau_{\rm el},-w,-x_K).
 \label{x_reciprocidad:rec:eq:reciproco-radial}
\end{equation}
Sur ce diamètre, les métriques hyperboliques de courbure \(-1\) satisfont
\begin{equation}
 \frac{|d\tau|^2}{(\Im\tau)^2}
 =\frac{4\,dw^2}{(1-w^2)^2}
 =\frac{dx_K^2}{(1-x_K^2)^2}
 =d\eta_{\rm el}^2=4\frac{dR_\star^2}{R_\star^2}.
 \label{x_reciprocidad:rec:eq:metricas}
\end{equation}
Ces égalités se vérifient en dérivant \eqref{x_reciprocidad:rec:eq:poincare}--\eqref{x_reciprocidad:rec:eq:klein}. En particulier, \(d(0,w)=|\eta_{\rm el}|\) et \(d(w,-w)=2|\eta_{\rm el}|\). La transformation \(w\mapsto-w\) est une rotation de \(\pi\) dans tout le disque et la symétrie par rapport au centre sur le diamètre engendré.

La normalisation \((Q/a_S,P/b_S)\) de l’intérieur de l’ellipse produit une autre réalisation du disque de Klein, avec des coordonnées de position. Sa métrique est
\[
 ds_K^2=\frac{(1-r^2)(du^2+dv^2)+(u\,du+v\,dv)^2}{(1-r^2)^2},
 \qquad r^2=u^2+v^2<1.
\]
Cette carte physique et la carte de forme marquée conservent des domaines différents. Dans la branche focale \(\eta_{\rm el}\ge0\), l'excentricité \(e_f=\sqrt{1-e^{-2\eta_{\rm el}}}\) donne la profondeur \(u_f=\operatorname{artanh}e_f\), d'où \(\eta_{\rm el}=\log\cosh u_f\). La profondeur focale n'est pas identifiée à l'anisotropie.

\begin{figure}[htbp]
\centering
\begin{tikzpicture}[>=Latex,scale=1.0,every node/.style={font=\small}]
\draw[ink,thick] (0,0) circle (1.55);
\draw[->] (-1.8,0)--(1.85,0);
\fill[accent] (-.75,0) circle (2pt) node[below=4pt] {\(w\)};
\fill[accent] (.75,0) circle (2pt) node[below=4pt] {\(-w\)};
\draw[<->,accent,thick] (-.75,.30)--(.75,.30);
\node at (0,-2.0) {Carte de forme : \(w\mapsto-w\)
};
\begin{scope}[xshift=7cm]
\draw[ink,thick] (0,0) circle (1.15);
\draw[->] (-1.6,0)--(2.05,0);
\draw[->] (0,-1.4)--(0,1.6);
\draw[accent,thick] (0,0)--(30:.65);
\draw[accent,thick,dashed] (30:.65)--(30:2.03);
\fill[accent] (30:.65) circle (2pt) node[above=3pt] {\(z\)};
\fill[accent] (30:2.03) circle (2pt) node[above=3pt] {\(1/\overline z\)};
\node at (0,-2.0) {Carte spectrale : inversion radiale
};
\end{scope}
\end{tikzpicture}
\caption{Deux réciprocités typées. La première reste à l’intérieur du disque et transporte l’orientation de la forme ; la seconde échange l’intérieur et l’extérieur du cercle unité. Les positions dessinées sont schématiques.}
\label{x_reciprocidad:fig:rec-dos-cartas}
\end{figure}

\section{Substitution structurelle des coordonnées angulaires}
\label{x_reciprocidad:rec:sustitucion}

La clôture dodécaphasique déjà construite publie la même coordonnée \(\alpha\) par
\[
 \alpha=\pi_{\HMT}+e_{\HMT}-\varphi_{\HMT}-4-\kappa_{\rm ag},
 \qquad \kappa_{\rm ag}=\frac{K_{\rm int}}{1000^{12}-1}.
\]
On maintient la base, l’origine de phase et la frontière de retenue compatibles. La récurrence analytique complète de l’article représente cette coordonnée corrélée ; sa publication n’est pas remplacée par un jet fini. En abrégeant les indices HMT, les lecteurs d’action de \eqref{x_reciprocidad:iv:action:h5}--\eqref{x_reciprocidad:iv:action:eta-ret} sont
\begin{align}
 H_5&=\frac12\sqrt{\frac{1000\alpha}{\varphi}}-\alpha
       +\frac9{16}\alpha^2-\frac59\alpha^3
       +\frac7{48}\alpha^4-\frac1{54}\alpha^5,\nonumber\\
 D_A&=e^{-100\pi\alpha/9},\qquad
 \mathcal C_\pi=169\alpha^6+\frac{D_A\alpha^7}{1-D_A\alpha},\nonumber\\
 \eta_{\rm ret}&=H_5-\frac{90}{\pi}\mathcal C_\pi,
 \qquad A=1000\alpha,\qquad C^*=2(\eta_{\rm ret}+\alpha).
 \label{x_reciprocidad:rec:eq:lectores-completos}
\end{align}
Le prolongement rationnel conserve tous les termes de la récurrence régionale. Sa convergence s'obtient sans troncature : pour \(\alpha>0\),
\[
 0<\alpha e^{-100\pi\alpha/9}\le\frac9{100\pi e}<1,
\]
car \(xe^{-cx}\) atteint son maximum \(1/(ce)\) en \(x=1/c\).

\begin{proposition}[Annulation structurelle dans le contraste angulaire]
\label{x_reciprocidad:rec:prop:sustitucion}
Le quotient engendré \(\xi=C^*/A\) satisfait exactement
\begin{align}
 \xi={}&\frac1{\sqrt{1000\varphi\alpha}}
 +\frac{9\alpha}{8000}-\frac{\alpha^2}{900}
 +\frac{7\alpha^3}{24000}-\frac{\alpha^4}{27000}\nonumber\\
 &-\frac9{50\pi}
 \left(169\alpha^5+\frac{D_A\alpha^6}{1-D_A\alpha}\right).
 \label{x_reciprocidad:rec:eq:xi-completa}
\end{align}
\end{proposition}
\begin{proof}
On part de \(\xi=(\eta_{\rm ret}+\alpha)/(500\alpha)\). Le terme \(-\alpha\) de \(H_5\) s’annule avec \(+\alpha\) du contre-angle. Diviser chaque terme restant par \(500\alpha\) produit l’expression : la racine donne \(1/\sqrt{1000\varphi\alpha}\), et \(90/500=9/50\). Le prolongement régional conserve son dénominateur complet.
\end{proof}

Par conséquent, \(x_K=-\xi\) et \(R_\star=((1-\xi)/(1+\xi))^{1/4}\) sont des expressions explicites des publications régionales et du registre. La substitution de la clôture de \(\alpha\) complète la dépendance de \(\pi,\varphi,e,\kappa_{\rm ag}\), sans promouvoir les reconnaissances métrologiques au rang de générateurs.

\section{Récupération de la coordonnée de Planck et conservation de l'orientation}

\begin{theorem}[Récupérateur radial de la section de retour]
\label{x_reciprocidad:rec:thm:accion}
Dans la carte positive, si \(q=R_\star^4\), alors
\begin{align}
 \eta_{\rm ret}&=\alpha\frac{499-501q}{1+q},\label{x_reciprocidad:rec:eq:eta-radio}\\
 \hbar_{\rm ret}&=\alpha\frac{499-501q}{1+q}\,10^{-34}\mathcal U_S,
 \qquad\hbar_{\rm pre}=H_5\,10^{-34}\mathcal U_S,\label{x_reciprocidad:rec:eq:planck-radio}\\
 R_{\rm act}&=\frac{\hbar_{\rm pre}}{\hbar_{\rm ret}}
 =\frac{H_5(1+q)}{\alpha(499-501q)}.\label{x_reciprocidad:rec:eq:ratio-radio}
\end{align}
La branche \(\eta_{\rm ret}>0\) équivaut à \(q<499/501\), à l’intérieur de \(q>0\).
\end{theorem}
\begin{proof}
Les identités angulaires donnent \(\eta_{\rm ret}=\alpha(500\xi-1)\). On substitue \(\xi=(1-q)/(1+q)\) ; le numérateur est \(499-501q\). Le dénominateur et \(\alpha\) sont positifs. La réalisation dimensionnelle et le quotient sont ceux des sections d'action préalablement construites. La décade \(-34\) et l'unité \(\mathcal U_S\) conservent leur dérivation déterminantale antérieure.
\end{proof}

Cette formule explicite la section de Planck dans le contre-angle. En termes de grandeurs d’action, la relation angulaire est
\[
 C^*=2\left(\frac{\hbar_{\rm ret}}{10^{-34}\mathcal U_S}+\alpha\right).
\]
La division par la base dimensionnelle extrait le coefficient adimensionnel ; elle n’ajoute pas une grandeur d’action à une constante sans unités.

Pour l’orientation angulaire \(\varepsilon\in\{-1,+1\}\), nous définissons
\[
 C^*_{\varepsilon}=2\varepsilon(\eta_{\rm ret}+\alpha),\quad
 \xi_\varepsilon=C^*_{\varepsilon}/A,\quad
 q_\varepsilon=\frac{1-\xi_\varepsilon}{1+\xi_\varepsilon}>0.
\]
Le récupérateur sur les deux cartes est
\begin{equation}
 \eta_{\rm ret}=\alpha\left(500\varepsilon
        \frac{1-q_\varepsilon}{1+q_\varepsilon}-1\right).
 \label{x_reciprocidad:rec:eq:eta-orientada}
\end{equation}

\begin{corollary}[Invariance du récupérateur conjugué]
La transformation \((\varepsilon,q)\mapsto(-\varepsilon,q^{-1})\) conserve \(\eta_{\rm ret}\), \(\hbar_{\rm ret}\) et \(R_{\rm act}\).
\end{corollary}
\begin{proof}
Le rapport \((1-q)/(1+q)\) change de signe lorsqu’on inverse \(q\). Le produit par \(\varepsilon\) reste invariant. On applique \eqref{x_reciprocidad:rec:eq:eta-orientada} puis les publications d’action.
\end{proof}

Si l’orientation est écartée et que \(\varepsilon=+1\) est indûment maintenu, la même expression produit \(\eta_{\rm ret}(q^{-1})=-\eta_{\rm ret}(q)-2\alpha\). Le registre distingue donc deux opérations algébriquement différentes. L’orientation angulaire \(\varepsilon\) ne s’identifie pas non plus sans autre précision au signe symplectique \(\sigma\) : l’échange des axes est antisymplectique, tandis que la rotation qui les échange préserve la forme d’aire.

\section{Entrelacement de phase et d'anisotropie}

Soient
\[
 S_\eta=\operatorname{diag}(e^{\eta/2},e^{-\eta/2}),\quad
 J_0=\begin{pmatrix}0&-1\\1&0\end{pmatrix},\quad
 J_\eta=S_\eta J_0S_\eta^{-1}
       =\begin{pmatrix}0&-e^\eta\\e^{-\eta}&0\end{pmatrix}.
\]
On a \(\det S_\eta=1\), \(J_\eta^2=-I\) et \(C_\eta(\theta)=e^{\theta J_\eta}=S_\eta e^{\theta J_0}S_\eta^{-1}\). La loi \(S_{\eta+\delta}=S_\delta S_\eta\) prouve l'entrelacement exact
\begin{equation}
 C_{\eta+\delta}(\theta)S_\delta=S_\delta C_\eta(\theta).
 \label{x_reciprocidad:rec:eq:entrelazamiento}
\end{equation}
La transformation préserve \(dQ\wedge dP\) et transporte l’évolution de phase vers la nouvelle ellipse. Avec la normalisation de \eqref{x_reciprocidad:rec:eq:elipse}, l’action balayée est \(\hbar\theta\), ou \(\sigma\hbar\theta\) dans le feuillet orienté. La phase \(\theta\), l’anisotropie \(\eta_{\rm el}\) et le paramètre \(t\) du flux de \(K\) conservent ainsi des fonctions différentes, reliées par des opérateurs explicites.

\section{Géodésique radiale et réalisation différentielle de la direction de K}

Chaque rayon de \eqref{x_reciprocidad:rec:eq:flujo} détermine une carte rectangulaire et une coordonnée hyperbolique :
\[
 \tau_i(t)=ir_i(t)^2,\qquad
 w_i(t)=\frac{r_i(t)^2-1}{r_i(t)^2+1}=\tanh(th_i).
\]

\begin{theorem}[Géodésique de rayons compensés]
\label{x_reciprocidad:rec:thm:geodesica}
Dans le produit de douze diamètres hyperboliques,
\begin{equation}
 d(w(s),w(t))=2|t-s|\,|\log\rho|,\qquad
 \sum_i\operatorname{artanh}w_i(t)=0.
 \label{x_reciprocidad:rec:eq:geodesica}
\end{equation}
La composition des avancées est
\[
 w_i(t+s)=\frac{w_i(t)+w_i(s)}{1+w_i(t)w_i(s)}.
\]
\end{theorem}
\begin{proof}
En coordonnées \(a_i=\log r_i\), la métrique produit est \(4\sum_i da_i^2\). La courbe \(a(t)=th\) est droite, de vitesse \(2\|h\|=2|\log\rho|\), et \(\sum_i a_i=0\). Chaque diamètre est totalement géodésique. La distance dans le produit est la racine de la somme des carrés des distances de ses facteurs, qui sont \(2|t-s||h_i|\). L'identité de somme nulle et la composition s'obtiennent à partir de \(\operatorname{artanh}w_i=th_i\) et de la formule d'addition de la tangente hyperbolique.
\end{proof}

Le rapport régional détermine la vitesse de cette réalisation ; \(K\) détermine sa direction. Si \(\rho=1\), la trajectoire est constante. Pour \(\rho\ne1\), l’espace radial de produit un est de dimension onze, de tangent \(V_{11}=\mathbf1^\perp\), et la normale à la trajectoire en son sein est
\[
 V_{10}(K)=V_{11}\cap u^\perp,\qquad
 P_{10}(K)=P_{11}-\frac{uu^{\mathsf T}}{\|u\|^2}.
\]
Le projecteur de la proposition~\ref{x_reciprocidad:prop:k-reduccion-diez} acquiert donc une réalisation différentielle : il élimine l’échelle uniforme et distingue la direction engendrée de ses dix variations transversales. Son domaine est cet espace radial, et non l’espace réel de dimension vingt-quatre de tous les disques.

Avec \(D_t=\operatorname{diag}(\operatorname{sech}^2(th_i))\), l’expression dans les coordonnées \(w\) est \(D_tP_{10}D_t^{-1}\). En effet, \(D_t\) est la différentielle de \(a\mapsto\tanh a\), et transforme la métrique \(4I\) en \(4\operatorname{diag}((1-w_i^2)^{-2})\). La projection transportée est orthogonale pour cette métrique.

\section{Relèvement réticulaire et relèvement rectangulaire}

Les rayons admettent deux réalisations qu’il convient de conserver explicites. L’opérateur \(G_t=\bigoplus_i r_iI_2\) dilate simultanément les deux axes de chaque plan \(A_2\). Il conserve le covolume total et commute avec l’isométrie ternaire. L’homothétie de chaque plan conserve son module complexe interne.

La carte \(\tau_i=ir_i^2\) possède une autre réalisation, d'aire fixe :
\[
 T_i=\operatorname{diag}(r_i^{-1},r_i),\qquad
 (\omega_1,\omega_2)=(r_i^{-1},ir_i),\qquad
 \omega_2/\omega_1=ir_i^2.
\]
C'est le relèvement elliptique avec \(\eta=-2\log r_i\). Cette convention transforme des périodes physiques. L'action homographique de la même matrice sur \(i\) donnerait \(i/r_i^2\) ; pour obtenir \(ir_i^2\) sous cette autre convention, on utilise \(T_i^{-1}\). La distinction des actions évite une inversion involontaire.

Les deux réalisations se reconstruisent à partir des mêmes rayons, mais la commutation ternaire de l’homothétie ne se transfère pas à la déformation anisotrope. L’état et ses lecteurs permettent de comparer les deux géométries sans identifier leurs opérateurs.
